\documentclass[11pt]{article}

\usepackage{amsmath, amssymb, geometry}
\usepackage{booktabs}
\usepackage{hyperref}

\geometry{margin=1in}

\title{Event Definition in GENEVA and Dependence on $r_N^{\mathrm{cut}}$}
\author{}
\date{}

\begin{document}

\maketitle


\textbf{Preserving \(r_{N-1}\)} in the projection/splitting map is not just a convenient choice. It is what makes the \textbf{\(N\)-jet spectrum differential in \(r_{N-1}\)} behave the way the GENEVA construction assumes at leading power. Let me make that precise.

\section*{Setup}

You are now considering the \textbf{\(N\)-event contribution above}
\[
r_{N-1}^{\rm cut},
\]
so the event is differential in an \(N\)-body phase-space point
\[
\Phi_N,
\]
and one of the coordinates of \(\Phi_N\) may be taken to be
\[
r_{N-1}.
\]

Then the \(N\)-event weight below \(r_N^{\rm cut}\) schematically contains
\[
\frac{d\sigma_N^{\rm MC}}{d\Phi_N}
\sim
\frac{d\sigma_N^{\rm scet}}{d\Phi_N}
+
\int d\Phi_{N+1}\,
\frac{d\sigma_{N+1}}{d\Phi_{N+1}}\,
\Theta(r_N<r_N^{\rm cut})\,
\delta\!\big(\Phi_N-\hat\Phi_N(\Phi_{N+1})\big)
\]
with support only for
\[
r_{N-1}(\Phi_N)>r_{N-1}^{\rm cut}.
\]

The question is: when \(\Phi_N\) is promoted to \(\Phi_{N+1}\) by an unresolved emission, must the map preserve \(r_{N-1}\)?

The answer is:

\subsection*{At leading power, yes.}

More precisely, the map should satisfy
\[
r_{N-1}\!\left(\hat\Phi_N(\Phi_{N+1})\right)
=
r_{N-1}^{\rm LP}(\Phi_{N+1})
\]
or at least differ from it only by subleading powers in \(r_N\).

In practice, the cleanest choice is to preserve it exactly.

\section*{Why \(r_{N-1}\) should be preserved}

\subsection*{1. Because \(r_{N-1}\) is now one of the measured Born variables}

Once your \(N\)-event distribution is differential in \(\Phi_N\), and \(\Phi_N\) includes \(r_{N-1}\), the projection delta is not merely assigning unresolved radiation to some generic \(N\)-body point. It is assigning it to a point with a \textbf{specific value of \(r_{N-1}\)}.

So if the map changes \(r_{N-1}\), then unresolved radiation generated from an \(N\)-body configuration with one value of \(r_{N-1}\) is deposited into a different bin in \(r_{N-1}\).

That immediately means the real-emission contribution below \(r_N^{\rm cut}\) is no longer local in the measured variable \(r_{N-1}\).

\subsection*{2. Because the singular  factorization/resummation is differential in that variable}

The singular  \(N\)-jet spectrum above \(r_{N-1}^{\rm cut}\) is organized as a distribution in \(\Phi_N\), hence also in \(r_{N-1}\).

The resummed spectrum and its FO expansion know about singular  structures of the form
\[
\frac{d\sigma^{\rm scet}}{d\Phi_N}
\equiv
\frac{d\sigma^{\rm scet}}{d r_{N-1}\, d\Omega_N}
\]
where \(d\Omega_N\) denotes the remaining \(N\)-body variables.

For the unresolved \((N+1)\)-body contribution below \(r_N^{\rm cut}\) to match this \textbf{point-by-point in \(r_{N-1}\)}, the projection must return the same \(r_{N-1}\) that appears in the singular  factorization formula, up to power corrections.

If not, then the cancellation/matching is no longer local in \(r_{N-1}\).

\subsection*{3. Because otherwise the delta expansion generates derivative terms in \(r_{N-1}\)}

Suppose the map does not preserve \(r_{N-1}\), and instead
\[
r_{N-1}(\hat\Phi_N)=r_{N-1}^{\rm scet}+\Delta r_{N-1},
\qquad
\Delta r_{N-1}\to0 \quad (r_N\to0).
\]

Then the projection delta contains
\[
\delta\!\big(r_{N-1}-r_{N-1}^{\rm scet}-\Delta r_{N-1}\big)
=
\delta\!\big(r_{N-1}-r_{N-1}^{\rm scet}\big)
-
\Delta r_{N-1}\,\partial_{r_{N-1}}
\delta\!\big(r_{N-1}-r_{N-1}^{\rm scet}\big)
+\cdots
\]

The first term is the leading-power singular  term, which is what the resummed-expanded spectrum knows how to reproduce.

The second term is not part of the LP factorization formula. After integrating the unresolved emission with its singular  kernel,
\[
\int_0^{r_N^{\rm cut}} \frac{dr_N}{r_N}\,\Delta r_{N-1},
\]
you get a residual contribution proportional to
\[
\partial_{r_{N-1}}\frac{d\sigma^{\rm scet}}{d\Phi_N}.
\]

So when \(r_{N-1}\) is not preserved, the mismatch is precisely a \textbf{migration term in the measured \(r_{N-1}\) spectrum}.

\section*{Physical meaning}

If you do not preserve \(r_{N-1}\), then unresolved radiation below \(r_N^{\rm cut}\) causes an event that should belong to one \(r_{N-1}\) bin to be reassigned to another one.

That has two consequences:

\subsection*{(a) Local spectrum distortion}

The total integrated \(N\)-event rate may still be correct up to power corrections, but the \textbf{differential spectrum in \(r_{N-1}\)} gets distorted by bin-to-bin migration.

\subsection*{(b) Loss of local cancellation}

The LP resummation/expansion cancels the singular  dependence locally at fixed \(r_{N-1}\). If the map shifts \(r_{N-1}\), this cancellation becomes nonlocal in \(r_{N-1}\), leaving a residual subleading term.

\section*{What happens parametrically when it is not preserved}

Let
\[
\Delta r_{N-1}\sim r_N^\alpha.
\]

Then the mismatch behaves like
\[
\int_0^{r_N^{\rm cut}} \frac{dr_N}{r_N}\,\Delta r_{N-1}
\sim
\int_0^{r_N^{\rm cut}} dr_N\, r_N^{\alpha-1}.
\]

So for the natural case \(\alpha=1\),
\[
\sim r_N^{\rm cut},
\]
and with logarithmic dependence in the coefficient one often gets
\[
\sim r_N^{\rm cut}\ln r_N^{\rm cut}.
\]

Thus the effect is formally a power correction in \(r_N^{\rm cut}\), but it directly multiplies derivatives of the \(r_{N-1}\) spectrum, so it can be numerically enhanced in regions where the spectrum is steep.

This is especially relevant:

\begin{itemize}
\item near cuts,
\item near peaks,
\item near Sudakov shoulders,
\item or wherever \(\frac{d\sigma}{d r_{N-1}}\) varies rapidly.
\end{itemize}

\section*{Why exact preservation is especially natural here}

There is a difference between a generic Born variable and \(r_{N-1}\) in your present setup.

Here \(r_{N-1}\) is not just some arbitrary coordinate of \(\Phi_N\). It is the \textbf{very variable that defines the \(N\)-jet event class above \(r_{N-1}^{\rm cut}\)} and often the variable whose spectrum you want to resum/match.

So preserving it ensures that:

\begin{itemize}
\item the unresolved \((N+1)\)-body contribution is assigned to the correct \(N\)-body point,
\item the matching to the singular  \(N\)-jet spectrum is local in the measured variable,
\item there is no spurious migration across \(r_{N-1}\) bins below \(r_N^{\rm cut}\).
\end{itemize}

\section*{What if it is not preserved?}

Then, even if the map is IR safe, you get:

\begin{itemize}
\item the same LP singular  structure after integrating over \(r_{N-1}\),
\item but a \textbf{different differential spectrum in \(r_{N-1}\)},
\item with residual terms of the form
\[
\sim
\partial_{r_{N-1}}
\left(\frac{d\sigma^{\rm scet}}{d\Phi_N}\right)
\times
\int_0^{r_N^{\rm cut}} \frac{dr_N}{r_N}\,\Delta r_{N-1}.
\]
\end{itemize}

In words:

\textbf{not preserving \(r_{N-1}\) induces subleading-power bin migration in the \(r_{N-1}\) spectrum.}

\section*{Very compact summary}

When the \(N\)-body phase space is differential in \(r_{N-1}\), the unresolved map
\[
\Phi_N\to\Phi_{N+1}\to\hat\Phi_N
\]
should preserve \(r_{N-1}\) so that the below-\(r_N^{\rm cut}\) contribution remains local in the measured variable \(r_{N-1}\).

If it does not, then unresolved emissions shift the event between \(r_{N-1}\) bins, and the difference shows up as a power-suppressed but potentially sizable distortion of the \(r_{N-1}\) spectrum, typically of order
\[
r_N^{\rm cut}
\quad\text{or}\quad
r_N^{\rm cut}\ln r_N^{\rm cut}.
\]

\section*{The sharpest statement}

For an \(N\)-jet spectrum differential in \(r_{N-1}\), preserving \(r_{N-1}\) is the analogue of preserving the Born variables used by the factorization theorem. It guarantees that the singular  FO contribution below \(r_N^{\rm cut}\) and the resummed/expanded \(N\)-jet spectrum are matched \textbf{locally in the measured spectrum variable}.

The cleanest way to see it is with a toy model for the \(N\)-jet spectrum differential in \(x \equiv r_{N-1}\).

\section*{Simple setup}

Assume the singular  unresolved \(N+1\)-body contribution below \(r_N^{\rm cut}\) has the form

\[
d\sigma_{N+1}^{\rm scet}
\;\sim\;
\frac{dr_N}{r_N}\, F(x)\,dx
\]

with \(x=r_{N-1}\), and the projection map assigns the unresolved configuration to

\[
x_{\rm proj} = x + \Delta x(r_N).
\]

Then the projected contribution to the \(x\)-spectrum is schematically

\[
\frac{d\sigma^{\rm proj}}{dx}
\sim
\int_0^{r_N^{\rm cut}} \frac{dr_N}{r_N}\,
F(x-\Delta x(r_N)).
\]

The resummed-expanded piece subtracts the leading-power term

\[
\int_0^{r_N^{\rm cut}} \frac{dr_N}{r_N}\,F(x),
\]

so the residual map-dependent mismatch is

\[
\Delta\!\left(\frac{d\sigma}{dx}\right)
\sim
\int_0^{r_N^{\rm cut}} \frac{dr_N}{r_N}
\Bigl[F(x-\Delta x(r_N)) - F(x)\Bigr].
\]

For small \(\Delta x\),

\[
F(x-\Delta x)-F(x)
\simeq
-\Delta x\,F'(x),
\]

hence

\[
\Delta\!\left(\frac{d\sigma}{dx}\right)
\sim
- F'(x)\int_0^{r_N^{\rm cut}} \frac{dr_N}{r_N}\,\Delta x(r_N).
\]

This is the master formula.

\section*{Case 1: \(r_{N-1}\) preserved exactly}

\[
\Delta x(r_N)=0
\]

so

\[
\Delta\!\left(\frac{d\sigma}{dx}\right)=0.
\]

No migration, no extra residual dependence from this source.

\section*{Case 2: \(r_{N-1}\) preserved only at leading power}

Suppose

\[
\Delta x(r_N) \sim c\,r_N.
\]

Then

\[
\Delta\!\left(\frac{d\sigma}{dx}\right)
\sim
- c\,F'(x)\int_0^{r_N^{\rm cut}} dr_N
=
- c\,F'(x)\,r_N^{\rm cut}.
\]

So the residual dependence is power suppressed:

\[
\sim r_N^{\rm cut} F'(x).
\]

This is the expected good behavior.

\section*{Case 3: \(r_{N-1}\) not preserved at leading power}

Now take a map for which the shift is larger, e.g.

\[
\Delta x(r_N)\sim c\,x
\qquad\text{or}\qquad
\Delta x(r_N)\sim c\,\sqrt{r_N}.
\]

These still vanish in the strict IR limit in the second example, but not fast enough.

\subsection*{Example A: \(\Delta x \sim c\,x\)}

Then

\[
\Delta\!\left(\frac{d\sigma}{dx}\right)
\sim
- c\,x\,F'(x)\int_0^{r_N^{\rm cut}} \frac{dr_N}{r_N}
=
- c\,x\,F'(x)\,\ln r_N^{\rm cut}.
\]

So instead of a power correction, you get a \textbf{logarithmic sensitivity to the cut}.

That is exactly the kind of behavior you do not want: the residual dependence is no longer suppressed as \(r_N^{\rm cut}\to 0\), but grows like \(\ln r_N^{\rm cut}\).

\subsection*{Example B: \(\Delta x \sim c\,\sqrt{r_N}\)}

Then

\[
\Delta\!\left(\frac{d\sigma}{dx}\right)
\sim
- c\,F'(x)\int_0^{r_N^{\rm cut}} \frac{dr_N}{\sqrt{r_N}}
=
-2c\,F'(x)\sqrt{r_N^{\rm cut}}.
\]

This still vanishes, but much more slowly than the desired \(r_N^{\rm cut}\). So it is parametrically enhanced compared to LP preservation.

\section*{Why this can blow up at small \(r_{N-1}\)}

Now take a steep physical spectrum near \(x=r_{N-1}=0\). For instance, a Sudakov-type behavior

\[
F(x)\sim \frac{1}{x}\,\exp[-L^2],\qquad L=\ln\frac{1}{x},
\]

or more simply, for illustration,

\[
F(x)\sim \frac{1}{x}.
\]

Then

\[
F'(x)\sim -\frac{1}{x^2}.
\]

Plug this into the bad case above.

\subsection*{If \(\Delta x\sim c\,x\)}

\[
\Delta\!\left(\frac{d\sigma}{dx}\right)
\sim
- c\,x \left(-\frac{1}{x^2}\right)\ln r_N^{\rm cut}
\sim
\frac{c}{x}\,\ln r_N^{\rm cut}.
\]

So as \(x=r_{N-1}\to 0\), the residual becomes large like

\[
\frac{\ln r_N^{\rm cut}}{r_{N-1}}.
\]

This is the sense in which the residual cut dependence can \textbf{blow up} in the physically important small-\(r_{N-1}\) region.

Even though the map is IR-safe in the strict \(r_N\to0\) sense, failure to preserve \(r_{N-1}\) at leading power creates a migration term multiplied by the steep derivative of the spectrum, and that becomes large precisely where resummation is supposed to give you the correct physical behavior.

\section*{Practical interpretation}

If \(r_{N-1}\) is not preserved at LP, then unresolved emissions below \(r_N^{\rm cut}\) shift the event by an amount that is too large compared to the power counting of the singular  measure \(dr_N/r_N\). The consequence is:

\begin{itemize}
\item the cancellation with the resummed-expanded result is no longer local in \(r_{N-1}\),
\item the residual dependence on \(r_N^{\rm cut}\) is no longer of the desired \(\mathcal O(r_N^{\rm cut})\) type,
\item and in the small-\(r_{N-1}\) region it can be amplified by the steep slope of the physical spectrum.
\end{itemize}

\section*{Compact takeaway to add}

A simple toy model makes the point transparent. Writing \(x=r_{N-1}\), the residual map dependence after FO–resummation matching behaves as

\[
\Delta\!\left(\frac{d\sigma}{dx}\right)
\sim
- F'(x)\int_0^{r_N^{\rm cut}}\frac{dr_N}{r_N}\,\Delta x(r_N).
\]

\begin{itemize}
\item If \(x\) is preserved exactly, \(\Delta x=0\), there is no such residual.
\item If \(x\) is preserved only at leading power, \(\Delta x\sim r_N\), the residual is \(\sim r_N^{\rm cut}F'(x)\), i.e. power suppressed.
\item If \(x\) is not preserved at leading power, e.g. \(\Delta x\sim x\), the residual becomes \(\sim xF'(x)\ln r_N^{\rm cut}\). For a steep spectrum \(F(x)\sim 1/x\), this gives \(\sim \ln r_N^{\rm cut}/x\), which becomes large as \(x=r_{N-1}\to0\).
\end{itemize}

This shows why preserving \(r_{N-1}\) at least at leading power is essential to obtain the correct physical small-\(r_{N-1}\) behavior once resummation is included.

\end{document}
    